\documentclass[border = 0mm]{standalone}

\usepackage[utf8]{inputenc}
\usepackage{amsmath}

\usepackage{xcolor}
% --
\usepackage{tikz}
\usetikzlibrary{scopes,angles,quotes,calc,intersections,math}

\usepackage{pgfplots}
\pgfplotsset{compat = 1.18}

\usepackage{fontspec}
\usepackage[default]{fontsetup}

\newcommand{\pointer}[1]{\node[red] at (#1) {$\times$};}

% ---
\begin{document}
\begin{tikzpicture}[
	line join = round, line cap = round,
]

	\tikzmath{
		\Width = 5.6cm; \Height = 3/4*\Width;
	};

	\node[
		inner sep = 0mm,
		minimum width = \Width,
		minimum height = \Height,	
	] (N) at (0,0) {};
	
	\clip (N.north west) rectangle (N.south east);
	\draw[line width = 0.5pt]
		([shift = {(2mm,-2mm)}]N.north west)
		rectangle
		([shift = {(-2mm,+2mm)}]N.south east);
				
	\node[align = center] at (0,0) {
		\large\bfseries
		Título \\[2mm]
		Autor 1\textsuperscript{a},\;
		Autor 2\textsuperscript{b} \\[2mm]
		\small\textsuperscript{a}
		\textit{Institución, ciudad.} \\
		\small\textsuperscript{b}
		\textit{Institución, ciudad.} \\[2mm]
		Año
	};

\end{tikzpicture}	
\end{document}